\SourcePage{104}
\section{一点辛几何}\label{sec:7}
\subsection{带角的正合辛流形}\label{sec:7a}
我们的基本设置涉及带角流形；相较于带边流形，其优点是这一类对乘积封闭，在讨论纤维化时很方便（见\sr{chapter:III}{第三章}）。先回顾有关术语。在带角流形中，考虑固定余维数 $k\geq1$ 的角点集合。这个集合的连通分支称为\emph{连通边界层}；这样的层在整个流形中的闭包称为\emph{连通边界面}。始终要求这些面没有自交；这意味着每个余维数为 $l$ 的角点，恰好属于余维数 $1\leq k\leq l$ 的 $2^l-1$ 个不同的连通边界面。去掉连通性要求，定义\emph{边界面}为同一余维数、两两不交的连通边界面的并；\emph{边界层}则为同一余维数、闭包两两不交的连通边界层的不交并。\znote{96}{$2^l-1$ 来自局部模型：在余维 $l$ 的角点附近，带角流形微分同胚于 $\{x\in\R^n:x^{n-l+1},\dots,x^n\geq0\}$ 中原点的一个邻域，此时恰有 $l$ 个坐标为零。过该点的每个连通边界层由「哪几个坐标为零、其余为正」指定，故有 $2^l-1$ 个非空选择，各层的闭包就是过该点的边界面；「无自交」的假设排除了更复杂的相遇方式。特别地，余维一的角点（$l=1$）恰属于一个边界面，这与带边流形的直觉一致。[来源：nLab, "manifold with boundary", https:/\allowbreak /\allowbreak ncatlab.org/\allowbreak nlab/\allowbreak show/\allowbreak manifold+\allowbreak with+\allowbreak boundary（2026-11-15 访问）。]}

\begin{definition*}
\emph{带角的正合辛流形} $(M,\varpi_M,\theta_M,I_M)$ 是紧致光滑带角流形 $M$，赋予辛形式 $\varpi_M$、满足 $d\theta_M=\varpi_M$ 的一形式 $\theta_M$，以及与 $\varpi_M$ 相容的殆复结构 $I_M$。要求满足以下两个凸性条件：由
\[
\varpi_M(\cdot,X_{\theta_M})=\theta_M
\]
定义的\emph{负 Liouville 向量场}沿 $M$ 的所有边界面都严格指向内部；边界还必须\emph{弱 $I_M$-凸}，即 $I_M$-全纯曲线不能接触 $\partial M$，除非完全包含在边界中。
\end{definition*}
这个定义是为兼顾各种需要而作的折中。缺点是两类凸性并未很好地联系起来，偶尔会使处理变得笨拙。当讨论只集中于一个流形时，省略 $\varpi_M,\theta_M,I_M$ 中的下标。

\begin{definition*}
带角正合辛流形的\emph{同构}是微分同胚 $\phi:M\to N$，满足 $\phi^*\theta_N=\theta_M+d(\text{某个函数})$，因而也满足 $\phi^*\varpi_N=\varpi_M$。
\end{definition*}
在所有这类同构构成的空间 $\operatorname{Iso}(M,N)$ 上赋予 $C^{\infty}$ 拓扑。$M=N$ 时，它成为拓扑群 $\operatorname{Aut}(M)$。注意，由于正合性假设，$\operatorname{Aut}(M)$ 中的任意同痕自动是 Hamilton 同痕。还将使用稍受限制、但技术上更简单的自同构概念：令 $\operatorname{Aut}(M,\partial M)$ 为在 $\partial M$ 附近等于恒等映射、且满足 $\phi^*\theta=\theta+d(\text{某个在 }\partial M\text{ 附近为零的函数})$ 的微分同胚 $\phi$ 所成的群。\znote{96}{这一句含一步验算。设 $\phi_t$ 是 $\operatorname{Aut}(M)$ 中的同痕，$\phi_t^*\theta=\theta+df_t$；对 $t$ 求导得 $\mathcal L_{X_t}\theta=d\dot f_t$，再用 Cartan 公式 $\mathcal L_{X_t}\theta=d(\theta(X_t))+\iota_{X_t}\varpi$，便得 $\iota_{X_t}\varpi=d(\dot f_t-\theta(X_t))$，故 $X_t$ 是 Hamilton 向量场，其 Hamilton 函数为 $\theta(X_t)-\dot f_t$。关键在右端自动恰当：$\varpi$ 自身正合，所以 $\iota_{X_t}\varpi$ 的 de Rham 类必为零，不存在 $H^1(M;\R)$ 型的障碍。反过来，$\varpi$ 正合也使每个 Hamilton 向量场的流都留在 $\operatorname{Aut}(M)$ 里：沿流 $\phi_t^*\varpi=\varpi$ 与 $\phi_t^*\theta=\theta+df_t$ 同时成立，$f_t$ 由 $\iota_{X_t}\varpi=dH_t$ 沿轨道积分给出。于是 $\operatorname{Aut}(M)$ 中的同痕与 Hamilton 同痕是同一批东西，$\operatorname{Iso}(M,N)$ 上取的 $C^\infty$ 拓扑也只反映 Hamilton 方向的变化。[来源：nLab, "Hamiltonian vector field", https:/\allowbreak /\allowbreak ncatlab.org/\allowbreak nlab/\allowbreak show/\allowbreak Hamiltonian+\allowbreak vector+\allowbreak field（2026-11-15 访问）。]}

\begin{definition*}
若 $\theta|_L$ 为正合一形式，就称 Lagrangian 子流形 $L\subset M$ \emph{正合}；除非另有说明，始终假设这些子流形是闭的，且与 $\partial M$ 不交。
\end{definition*}
这一类子流形在 $\operatorname{Aut}(M)$ 中的同痕以及负 Liouville 向量场的流下保持不变。\znote{96}{正合性的作用在此：$\varpi|_L=d(\theta|_L)$ 恒成立，故 $\dim L=\tfrac12\dim M$ 且 $\theta|_L$ 闭的子流形已是 Lagrangian；要求 $\theta|_L$ 整体恰当，即 $[\theta|_L]=0\in H^1(L;\R)$，是更强的条件，等价于 $L$ 上任何闭曲线的辛面积为零。有了整体原函数 $h_L$，$\phi^*\theta=\theta+df$ 直接把正合性搬过去：$\phi^*(\theta|_{\phi(L)})=d(h_L+f|_L)$，故同痕像仍正合；负 Liouville 流的情形同理。这正是这一类子流形能作为本章「对象」的前提：$\theta|_L$ 的整体原函数 $h_L$ 使 $L$ 上闭曲线的辛面积为零，作用量泛函因而无需 Novikov 系数即可定义。[来源：nLab, "lagrangian submanifold", https:/\allowbreak /\allowbreak ncatlab.org/\allowbreak nlab/\allowbreak show/\allowbreak Lagrangian+\allowbreak submanifold（2026-11-15 访问）。]}

\begin{definition*}
带角正合辛流形的\emph{变形}，就是固定流形 $M$ 上的一参数数据族 $(\varpi^s,\theta^s,I^s)$（$s\in[0,1]$），对每个 $s$ 都
\SourcePage{105}
满足上述条件。
\end{definition*}
这一背景下有 Moser 引理的一个较弱对应物，可表述如下。设 $M,N$ 是带角正合辛流形。
\begin{definition*}
若对某个 $c\in\R$ 有 $\rho^*\theta_N=e^{-c}\theta_M+d(\text{某个函数})$，就称嵌入 $\rho:M\to N$ 为\emph{正合共形辛嵌入}。
\end{definition*}
\begin{lemma}\label{st:7.1}
若两个带角正合辛流形可以相互变形，则存在从其中一个到另一个的正合共形辛嵌入。
\end{lemma}
\begin{proof}
取一个变形，像标准 Moser 论证那样，考察由 $\partial\theta^s/\partial s$ 得到的依赖于 $s$ 的向量场。它未必沿边界面指向内部，但可以通过加上 $X_{\theta^s}$ 的一个足够大的倍数来补救。积分所得向量场，即得所需嵌入。
\end{proof}
还可以说得更精确些。在所有正合共形辛嵌入的空间中，考察刚才构造的嵌入 $\rho$ 的同痕类。它不依赖于任何选择；在保持端点不变的变形同伦下不变；并且与变形的复合相容。特别地，若 $\tilde\rho$ 是将变形反向进行所得到的嵌入，那么 $\tilde\rho\circ\rho$ 与恒等映射同痕。

\subsection{接触型边界的情形}\label{sec:7b}
为补充上面的讨论，现在考虑一类更熟悉的带边辛流形。一些基础参考文献为\bcite{42}、\bcite{100}。
\begin{definition*}
\emph{具有接触型边界的正合辛流形}是一个紧致带边而不带角的流形 $M$，赋予辛形式 $\varpi_M$ 与一形式 $\theta_M$，满足 $d\theta_M=\varpi_M$，且 $X_{\theta_M}$ 沿 $\partial M$ 指向内部。
\end{definition*}
一个基本构造，是以典范方式接上一个无限锥，将这样的 $M$ 扩大为非紧致正合辛流形：
\begin{equation}\label{eq:7.1}
\begin{aligned}
(\widehat M,\varpi_{\widehat M},\theta_{\widehat M})
={}&(M,\varpi_M,\theta_M)\cup_{\partial M}\\[-2pt]
&\bigl(\R^+\times\partial M,\ d(e^r(\theta_M|_{\partial M})),\ e^r(\theta_M|_{\partial M})\bigr),
\end{aligned}
\end{equation}
其中 $r$ 为锥的径向变量；全书约定 $\R^{\pm}\subset\R$ 是闭的，即都包含零。“接缝”处的可微结构，利用两边 Liouville 流给出的 $\partial M$ 的领状邻域来定义。因此，$M$ 上的负 Liouville 场 $X_{\theta_M}$ 与 $\R^+\times\partial M$ 上的 $-\partial_r$ 拼合为 $\widehat M$ 上的负 Liouville 场。类似地，在 $\partial M$ 的某个邻域 $U\subset M$ 中，存在唯一的函数芽 $h_M:U\to(0,1]$，满足
\begin{equation}\label{eq:7.2}
h_M^{-1}(1)=\partial M,\qquad X_{\theta_M}\cdot h_M=-h_M.
\end{equation}
在锥上令 $h_{\widehat M}(r,y)=e^r$，就可将它光滑延伸为 $U\cup(\R^+\times\partial M)\subset\widehat M$ 上的函数 $h_{\widehat M}$。\znote{97}{$h_M$ 可读作「到边界的 Liouville 距离」：$X_{\theta_M}\cdot h_M=-h_M$ 说明沿负 Liouville 流 $\phi_t$（$\dot\phi_t=X_{\theta_M}$）有 $h_M(\phi_t x)=e^{-t}h_M(x)$，而 $h_M^{-1}(1)=\partial M$ 把它归一化。正因如此，锥上取 $h_{\widehat M}=e^r$ 才与 $\theta_{\widehat M}=e^r(\theta_M|_{\partial M})$ 配对，负 Liouville 场拼成 $-\partial_r$。芽的唯一性也随之清楚：$h_M$ 沿每条轨道被端点值唯一确定，故它在 $\partial M$ 的邻域内没有选择余地。顺带一提，后面引理 7.5 证明中取的 $M^{\epsilon}=M\setminus h_M^{-1}((e^{-\epsilon},1])$ 正是用 $h_M$ 把边界推里一点，而 $h_M$ 在 $1$ 附近的水平集给出 $\partial M$ 的领状邻域。[来源：原书 (7.1)--(7.2) 及本节关于 $h_{\widehat M}$ 延拓的论述。]}

\SourcePage{106}
\begin{definition*}
两个\eref{7.1}类型的流形之间，适当的同构概念是微分同胚 $\hat\phi:\widehat M\to\widehat N$，满足 $\hat\phi^*\theta_{\widehat N}-\theta_{\widehat M}=d(\text{某个紧支集函数})$。
\end{definition*}
特别地，这样的映射在某个紧子集之外保持 Liouville 向量场 $\partial_r$。因此，对某个 $C\gg0$，它在 $[C,\infty)\times\partial M\subset\widehat M$ 上形如
\[
\hat\phi(r,y)=(r-w(y),\psi(y)),
\]
其中 $\psi:\partial M\to\partial N$ 为接触微分同胚，$w\in C^{\infty}(\partial M,(-\infty,C])$ 满足 $\psi^*(\theta_N|_{\partial N})=e^w(\theta_M|_{\partial M})$。称 $\hat\phi$ \emph{在无穷远附近为接触型}。所有这样的同构构成的空间 $\operatorname{Iso}(\widehat M,\widehat N)$ 带有自然拓扑；它允许同痕改变无穷远处的接触微分同胚，因此这些同痕未必具有紧支集。照常记 $\operatorname{Aut}(\widehat M)=\operatorname{Iso}(\widehat M,\widehat M)$。以下 Moser 引理的版本是标准事实：
\begin{lemma}\label{st:7.2}
设 $(\varpi_M^s,\theta_M^s)$（$s\in[0,1]$）是 $M$ 上具有接触型边界的正合辛结构的光滑族。那么，由此得到的非紧致流形 $\widehat M^0,\widehat M^1$ 之间，存在在无穷远附近为接触型的正合辛同构。
\end{lemma}
设 $M$ 为具有接触型边界的正合辛流形。沿用一个众所周知的方法，例如见\bcite{76}，考虑与 $\varpi_M$ 相容的殆复结构 $J$，要求它在 $\partial M$ 的某个邻域中在 $X_{\theta_M}$ 的流下不变，并满足
\begin{equation}\label{eq:7.3}
dh_M\circ J=-\theta_M,
\end{equation}
其中 $h_M$ 如\eref{7.2}。称这样的殆复结构\emph{在边界附近为接触型}。给定带复结构 $I_S$ 的 Riemann 曲面 $S$ 与 $J$-全纯映射 $u:S\to M$，考虑函数 $\rho=h_M\circ u:u^{-1}(U)\to(0,1]$。由\eref{7.3}，它满足
\begin{equation}\label{eq:7.4}
\begin{aligned}
\Delta\rho&=-d(d\rho\circ I_S)=-d(dh_M\circ J\circ du)\\
&=d(\theta_M\circ du)=u^*\varpi_M\geq0.
\end{aligned}
\end{equation}
由极大值原理，若 $\rho$ 在 $S$ 的某个内点达到其最大可能值一，则它在那里必然局部为常数。因此，任取在边界附近为接触型的 $I_M=J$，就能使 $\partial M$ 弱凸，从而使 $M$ 成为带角正合辛流形。

我们需要这一基本凸性论证的几个变体。首先，它在以下意义下适用于殆复结构族：
\begin{lemma}\label{st:7.3}
设 $(J^z)_{z\in S}$ 是 $M$ 上与 $\varpi_M$ 相容的殆复结构族，以某个 Riemann 曲面 $S$ 为参数。假设每个 $J^z$ 都在 $\partial M$ 附近为接触型。那么，$S\times M$ 上由 $J(z,x)=I_S(z)\times J^z(x)$ 给出的殆复结构使 $S\times\partial M$ 弱凸。
\end{lemma}
\SourcePage{107}
其次，可以考察完备化\eref{7.1}中的伪全纯曲线。$M$ 上每个在边界附近为接触型的殆复结构，都有到 $\widehat M$ 的典范延伸 $\hat J$，在锥上对于 $r$ 的平移不变。它满足 $dh_{\widehat M}\circ\hat J=-\theta_{\widehat M}$，因此用与前面相同的计算，得到：
\begin{lemma}\label{st:7.4}
设 $S$ 是紧致连通带边 Riemann 曲面，$u:S\to\widehat M$ 是 $\hat J$-全纯映射，满足 $u(\partial S)\subset M$。那么也有 $u(S)\subset M$。
\end{lemma}
极大值原理论证还有积分形式；它既可以给出前一引理的另一种证明，也能得到稍为一般的结果。具体地，设 $N$ 为任意正合辛流形，$M$ 为具有接触型边界的同维正合辛流形，$\iota:M\hookrightarrow\operatorname{int}(N)$ 为满足 $\omega_M=\iota^*\omega_N$、$\theta_M=\iota^*\theta_N$ 的嵌入。为简洁起见，把它写作包含。取与 $\omega_N$ 相容的殆复结构 $J$，其在 $M$ 上的限制在边界附近为接触型。令 $S$ 为紧致连通带边 Riemann 曲面，$u:S\to N\setminus\operatorname{int}(M)$ 为满足 $u(\partial S)\subset\partial M$ 的 $J$-全纯映射。显然，
\begin{equation}\label{eq:7.5}
\begin{aligned}
0\leq\int_Su^*\omega_N&=\int_{\partial S}u^*\theta_M=\int_{\partial S}u^*(-dh_M\circ J)\\
&=\int_{\partial S}-dh_M\circ du\circ I_S.
\end{aligned}
\end{equation}
若 $\xi$ 是沿 $\partial S$、对于边界定向为正的切向量，则 $I_S\xi$ 指向内部，而 $du(I_S\xi)$ 不指向 $M$ 的内部，因此 $dh_M(du(I_S\xi))\geq0$。除非 $u$ 为常值，否则这与\eref{7.5}矛盾。作为一个简单应用，得到：
\begin{lemma}\label{st:7.5}
设 $S$ 是紧致连通带边 Riemann 曲面，$u:S\to N$ 是 $J$-全纯映射，满足 $u(\partial S)\subset\operatorname{int}(M)$。那么也有 $u(S)\subset\operatorname{int}(M)$。
\end{lemma}
\begin{proof}
对于某个 $\epsilon>0$，以稍小的子集 $M^{\epsilon}=M\setminus h_M^{-1}((e^{-\epsilon},1])$ 替换 $M$。适当选择 $\epsilon$，可以使 $J|_{M^{\epsilon}}$ 在边界附近仍为接触型，使 $u(\partial S)$ 位于 $\operatorname{int}(M^{\epsilon})$ 中，并使 $u$ 与 $\partial M^{\epsilon}$ 横截。因此 $S^{\epsilon}=u^{-1}(N\setminus\operatorname{int}(M^{\epsilon}))$ 仍为紧致带边 Riemann 曲面，其边界映到 $\partial M^{\epsilon}$。对其每个连通分支应用前面的论证，得到 $u(S^{\epsilon})\subset M^{\epsilon}$，从而 $u(S)\subset\operatorname{int}(M)$。
\end{proof}
最后，给定任意带角正合辛流形 $N$，可以如下\emph{磨圆角点}。令 $F_1,\ldots,F_d$ 为 $N$ 的余维一边界面。对每个 $F_j$，选择光滑函数 $f_j:N\to\R^+$，满足 $F_j=f_j^{-1}(0)$，且在 $F_j$ 的向外法向量上 $df_j<0$，从而特别有
\SourcePage{108}
$X_{\theta_N}\cdot f_j>0$ 在 $F_j$ 的某个邻域中成立。令 $f=f_1f_2\cdots f_d$。我们断言，在 $\operatorname{int}(N)$ 中充分靠近边界的所有点处，$X_{\theta_N}\cdot f>0$。这一不等式可以改写为
\begin{equation}\label{eq:7.6}
\sum_j\frac{X_{\theta_N}\cdot f_j}{f_j}>0.
\end{equation}
设 $\operatorname{int}(N)$ 中的一列点收敛到某个极限，该极限属于 $F_j$ 当且仅当 $j\in E$，其中 $E\subset\{1,\ldots,d\}$ 为某个非空子集。沿此点列，\eref{7.6}中 $j\in E$ 的各项都趋于 $+\infty$，其余各项有界，故该不等式最终成立。证明了这一断言后，立即可见：
\begin{lemma}\label{st:7.6}
对任意充分小的正数 $\delta$，$M=f^{-1}([\delta,\infty))\subset N$ 配备 $\varpi_N,\theta_N$ 的限制后，是具有接触型边界的正合辛流形。此外，经过足够长的时间，负 Liouville 流将把整个 $N$ 压入 $\operatorname{int}(M)$。
\end{lemma}
